\begin{afterword}
\summaryhead

The post starts from old science-fiction magazine covers on which alien monsters carry off women in torn dresses. A non-human alien with a different evolutionary history would have no reason to find a human woman attractive. The author suggests that the artists made this mistake not by reasoning that all minds are alike, but because they perceived sexiness as a property of the woman herself, like her height, rather than as a response in the mind of whoever looks at her.

The author names this error, following the physicist E. T. Jaynes, the Mind Projection Fallacy: treating a property of one's own mind as a property of the external world. Jaynes applied it mainly to probability, which he treated as a state of partial knowledge rather than a feature of objects; the author promises to return to this. The post lists other cases, including the dispute over whether an unheard falling tree makes a sound, Kant on space and Hume on a priori knowledge, and ends with two humorous examples from science fiction.

\responsehead

This post names an error, gives one example of it, and stops. The example does not show the error, the list that generalizes it names two philosophers and quotes neither accurately, and the argument about probability that the name was introduced for is deferred to another post.

The example is a pulp cover, and the post's diagnosis of it is a guess (``Probably the artist did not even think to ask''). An obvious alternative is not considered. A cover is there to sell the magazine, so the woman may be on it for the buyer, not for the monster. On that reading the illustrator believed nothing about aliens at all. From this one case the post moves to a claim about how ``we instinctively represent'' attraction, and then hedges the model until it forbids nothing: ``that data structure, or something metaphorically similar to it.''

The generalization is worse. Kant is listed for declaring space ``flat.'' Kant wrote that ``Space represents no property at all of any things in themselves''; for him space is a form of our intuition. He located space in the mind; the post files him under projecting the mind onto the world. Hume is listed for a sentence that says certain propositions hold ``without dependence on what is any where existent,'' misquoted and misdescribed as a definition of ``a priori ideas.'' And Hume is the philosopher who described projection: ``the mind has a great propensity to spread itself on external objects.'' The post credits Jaynes, a twentieth-century physicist, with the name, lists the eighteenth-century writer who had the idea as an example of the error, and does not mention what he wrote.

The post has no conclusion. After the list it ends on two anecdotes about romance with aliens and the line ``Just thought I'd mention that.''

In short: a borrowed name, one example with a simpler explanation, and a list of philosophers that files Kant, who put space in the mind, under projection, and enrolls Hume, the source of the idea, as an offender.
\end{afterword}
